%%%%%%%%%%%%%%%%%%%%%%%%%%%%%%%%%%%%%%%%%%%%%%%%%%%%%%%%%%%%%%%%%%%%%%%%%%%%%%%%
\section{Automatic Software Deployment}
{   
	\usebackgroundtemplate{
		\vbox to \paperheight{\vfil\hbox to \paperwidth{\hfil\includegraphics[height=\paperheight]{misc/deployment.png}\hfil}\vfil}
		% source: https://en.m.wikipedia.org/wiki/File:Text-x-python.svg
	}
	\frame{
		\frametitle{Software Deployment with \Snakemake}
		\begin{mdframed}[tikzsetting={draw=white,fill=white,fill opacity=0.8,
				line width=0pt},backgroundcolor=none,leftmargin=0,
			rightmargin=150,innertopmargin=4pt,roundcorner=10pt]
			\tableofcontents[currentsection,sections={1-4},hideothersubsections]
		\end{mdframed}
	}
}

%%%%%%%%%%%%%%%%%%%%%%%%%%%%%%%%%%%%%%%%%%%%%%%%%%%%%%%%%%%%%%%%%%%%%%%%%%%%%%%%
\begin{frame}
	\frametitle{What is this about?}
	\begin{question}[Questions]
		\begin{itemize}
			\item Which software deployment methods are supported in \Snakemake?
			\item How can I use them as a user / as a developer?
		\end{itemize}
	\end{question}
	\begin{docs}[Objectives]
		\begin{enumerate}
			\item Learn to use the supported methods within \Snakemake for software deploymend
			\item Know how to support \Snakemake's default method: Conda.
		\end{enumerate}
	\end{docs}
\end{frame}

%%%%%%%%%%%%%%%%%%%%%%%%%%%%%%%%%%%%%%%%%%%%%%%%%%%%%%%%%%%%%%%%%%%%%%%%%%%%%%%%
\begin{frame}
	\frametitle{Basics of Software Deployment}
	\Snakemake provides a \altverb{--software-deployment-flag} or short \altverb{--sdm}. It supports three modules:
	\begin{itemize}
		\item Conda
		\item Container (Apptainer)
		\item Environment Modules
	\end{itemize}
    \begin{hint}
    	It is possible to use them in a mix for a given workflow.
    \end{hint}
\end{frame}

%%%%%%%%%%%%%%%%%%%%%%%%%%%%%%%%%%%%%%%%%%%%%%%%%%%%%%%%%%%%%%%%%%%%%%%%%%%%%%%%
\begin{frame}[fragile]
	\frametitle{Deployment for Conda in Practice}
	  We can instruct Snakemake to use Conda:
	  \begin{lstlisting}[language=Bash, style=Shell]
$ snakemake --sdm conda 
	\end{lstlisting}
    \begin{columns}
    	\begin{column}{0.5\textwidth}
    		This requires rules to specify the a Conda environment:
    		\begin{lstlisting}[language=Python,style=Python]
rule NAME:
    input:
    	"table.txt"
    output:
    	"plots/myplot.pdf"
    conda:
    	"envs/ggplot.yaml"
    script:
        "scripts/plot-stuff.R"
    		\end{lstlisting}
    	\end{column}
        \begin{column}{0.5\textwidth}
        	\pause
        	with an environment definition like this:
            \begin{lstlisting}[language=Python,style=Python]
channels:
- conda-forge
- bioconda
dependencies:
- r=3.3.1
- r-ggplot2=2.1.0
           \end{lstlisting}	
        \end{column}
    \end{columns}
\end{frame}

%%%%%%%%%%%%%%%%%%%%%%%%%%%%%%%%%%%%%%%%%%%%%%%%%%%%%%%%%%%%%%%%%%%%%%%%%%%%%%%%
\begin{frame}[fragile]
	\frametitle{Deployment for Conda in Practice - II}
	\begin{warning}
		Hic sunt dracones: Just using \altverb{--sdm} might lead to trouble!
	\end{warning}
	\begin{columns}
		\begin{column}{0.5\textwidth}
			This requires rules to specify the a Conda environment:
			\begin{lstlisting}[language=Bash, style=Shell]
$ # deploy with conda
$ snakemake --sdm conda \ 
> \ # avoid overloading HOME
> --conda-prefix /path/to/project/dir \ 
> \ # cleanup pckgs and tarballs
> --conda-cleanup-pkgs 
			\end{lstlisting}
		\end{column}
		\begin{column}{0.5\textwidth}
			\pause
			When a workflow aborts due to quota limitations, you can perform:
			\begin{lstlisting}[language=Bash,style=Shell]
(my_env) $ # clean up all pckgs 
(my_env) $ # and tarballs
(my_env) $ # in a given environment
(my_env) $ conda clean -a 
			\end{lstlisting}
		    \begin{docs}
		       On multiuser systems the concern is to use too many "inodes" ($\propto$ the number of files), hence most systems have strict quotas.
		    \end{docs}
		\end{column}
	\end{columns}
\end{frame}

%%%%%%%%%%%%%%%%%%%%%%%%%%%%%%%%%%%%%%%%%%%%%%%%%%%%%%%%%%%%%%%%%%%%%%%%%%%%%%%%
\begin{frame}[fragile]
	\frametitle{Updating Environments}
	A workflow with many dependencies, if not maintained can be outdated. \altverb{snakedeploy} has an auto-update feature:
	\begin{lstlisting}[language=Bash,style=Shell]
$ snakedeploy update-conda-envs envs/ggplot.yaml
    \end{lstlisting}
    You are also able to update multiple environments:
    \begin{lstlisting}[language=Bash,style=Shell]
$ snakedeploy update-conda-envs envs/*yaml
    \end{lstlisting}
    \pause
    \begin{warning}
      Occasionally, \altverb{--conda-frontend mamba} can be more stable.
    \end{warning}
\end{frame}

%%%%%%%%%%%%%%%%%%%%%%%%%%%%%%%%%%%%%%%%%%%%%%%%%%%%%%%%%%%%%%%%%%%%%%%%%%%%%%%%
\begin{frame}[fragile]
	\frametitle{Using Conntainers}
	\begin{columns}
		\begin{column}{0.5\textwidth}
	        Our ggplot-example can also be served by a container:
			\begin{lstlisting}[language=Python,style=Python]
rule NAME:
    input:
        "table.txt"
    output:
        "plots/myplot.pdf"
    container:
       "docker://joseespinosa/docker-r-ggplot2:1.0"
    script:
       "scripts/plot-stuff.R"
			\end{lstlisting}
		\end{column}
		\begin{column}{0.5\textwidth}
			\pause
			For deployment use
			\begin{lstlisting}[language=Bash,style=shell]
snakemake --sdm apptainer
			\end{lstlisting}
		    \begin{itemize}
		    	\item Apptainer is used for security reasons.
		    	\item Allowed image urls entail everything supported by apptainer (e.g., \altverb{shub://} and \altverb{docker://}). However, \altverb{docker://} is preferred.
		    \end{itemize}	
		\end{column}
	\end{columns}
\end{frame}

%%%%%%%%%%%%%%%%%%%%%%%%%%%%%%%%%%%%%%%%%%%%%%%%%%%%%%%%%%%%%%%%%%%%%%%%%%%%%%%%
\begin{frame}[fragile]
	\frametitle{Using Environment Modules}
	\begin{columns}
		\begin{column}{0.5\textwidth}
			Our mapping rule could be served with:
			\begin{lstlisting}[language=Python,style=Python]
rule bwa:
...
    conda:
        "envs/bwa.yaml"
    envmodules:
        "bio/bwa/0.7.9",
        "bio/samtools/1.9"
    shell:
        "bwa mem {input} | \"
        " samtools view -Sbh - > {output}"
			\end{lstlisting}
		\end{column}
		\begin{column}{0.5\textwidth}
			\pause
			For deployment use
			\begin{lstlisting}[language=Bash,style=shell]
snakemake --sdm envmodules
			\end{lstlisting}	
		\end{column}
	\end{columns}
\end{frame}
